\documentclass{eceasst}
% This is the source of the author documentation
% for the ECEASST document class.

% Required packages
% =================
\usepackage{subfig}

\usepackage{orcidlink}
\usepackage{accsupp}
\usepackage{todonotes}

%Mathematical formula accessibility and machine-readability: if you add
\usepackage[accsupp]{axessibility}
\pdfcompresslevel=1
%in the article preamble, you can select an equation and copy-paste it to a text file to recover the TeX math


% PDF accessibility and PDF metadata
\newcommand{\authorRef}[1]{\texorpdfstring{\autref{#1}}{}}
\newcommand{\authorOrcid}[1]{\texorpdfstring{\thinspace\orcidlink{#1}\thinspace}{}}
\newcommand{\geprislink}[1]{\href{https://gepris.dfg.de/gepris/projekt/#1?language=en}{#1}}

\newcommand{\Eg}{E.\,g.}
\newcommand{\eg}{e.\,g.}

% ================================
% Environment for AI Use
% ================================
\newenvironment{aiuse}{%
  \par\vspace{11pt}%
  \noindent\textbf{Declaration on the use of AI:}\enspace}{%
  \par\vspace{11pt}}

% ================================
% Environment for Data availability
% ================================
\newenvironment{data}{%
  \par\vspace{11pt}%
  \noindent\textbf{Data availability:}\enspace}{%
  \par\vspace{11pt}}



\definecolor{easstblue}{rgb}{.05,.32,.66}



\title{INF Projects in Germany's CRCs --- an overview and lessons learned}
\short{INF Projects in Germany's CRCs}

\author{
Hamza Oukili\authorOrcid{0000-0003-4695-8532}\authorRef{1},
Florian Goth\authorOrcid{0000-0003-2707-4790}\authorRef{2},
Frank Löffler\authorOrcid{0000-0001-6643-6323}\authorRef{3},
Mercedes Martinez Bruera\authorOrcid{0009-0004-2891-7045}\authorRef{4},
Ron Dockhorn\authorOrcid{0000-0002-5268-5430}\authorRef{8},
Alvaro Aguilera\authorOrcid{0000-0001-9324-8555}\authorRef{8},
Fabian~Prasser\authorOrcid{0000-0003-3172-3095}\authorRef{5},
Michael~Gnädig\authorOrcid{}\authorRef{5},
Marco~Johns\authorOrcid{0000-0003-4706-6595}\authorRef{5},
Maija~Poikela\authorOrcid{0000-0003-2616-4401}\authorRef{5}, 
Tobias~Penzkofer\authorOrcid{0000-0001-9591-8575}\authorRef{6}, 
Lisa~Franziska~Buchauer\authorOrcid{0000-0002-4722-8390}\authorRef{6}, 
Aristotelis~Misios\authorOrcid{}\authorRef{6}
Maryam Mohammadi\authorOrcid{0009-0007-8747-029X}\authorRef{7}
Lincoln Sherpa\authorOrcid{0009-0000-3287-0295}\authorRef{8}
} % Authors and references to addresses

\institute{% Institutes with labels
\autlabel{1} Institute for Modeling Hydraulic and Environmental Systems (IWS), University of Stuttgart,\\
\autlabel{2} Institut für theoretische Physik 1, University of Würzburg, 97074, Würzburg,\\
\autlabel{3} Heinz Nixdorf Chair for Distributed Information Systems, Jena, Germany,\\
\autlabel{4} Goethe-Universität Frankfurt am Main, Frankfurt am Main,\\
\autlabel{8} Center for Interdisciplinary Digital Sciences, Technische Universität Dresden, D-01069 Dresden, Germany,\\
\autlabel{5} Berlin Institute of Health at Charité - Universitätsmedizin Berlin, Berlin, Germany\\
\autlabel{6} Charité - Universitätsmedizin Berlin, Berlin, Germany\\
\autlabel{7} Linguistic Department, University of Bielefeld\\
}

\abstract{% Abstract of the article
In the Research funding landscape of Germany, so-called collaborative Research Centers (CRCs) play a crucial role in the competitive funding steered by the German Research Foundation.
Digital tools have become so common in academic research that more and more processes in the research cycle require their use.
It has become commonplace to pool certain tasks in dedicated units within the CRCs, the so-called INF projects.
This report tries to catalogue the INF projects in Germany's research landscape.
Building on that and on the workshop held at the de-RSE26 conference we collect lessons learned across INF projects and distill the discussion into guides that other INF projects (or similar structures) can follow.
}
\keywords{INF, CRCs, DFG, Research Data Management, RDM, Research Software Engineering, RSE}



% PDF/A compliance
% ========================
% This is an optional step to help with archiving and long-term preservation
% of electronic documents. It stores XMP metadata, machine-readable content,
% and color profiles in the PDF file. When used, it must appear last in the
% preamble to avoid conflicts with other packages (the pdfx package loads
% hyperref and xcolor).

\usepackage[a-3u]{pdfx}
\begin{filecontents}[overwrite]{\jobname.xmpdata}
\Author{Hamza Oukili, Florian Goth, Bernd Flemisch et al.}
\Title{INF Projects in Germany's CRCs --- an overview and lessons learned}
\Language{English}
\Subject{Software Science and Technology}
\end{filecontents}


\begin{document}
\maketitle

\section{Introduction}
% TODO: make this different from the abstract
In the Research funding landscape of Germany, so-called collaborative Research Centers (CRCs) play a crucial role in the competitive funding steered by the German Research Society.
Digital tools have become so common in academic research that more and more processes in the research cycle require their use.
It has become commonplace to pool certain tasks in dedicated units within the CRCs, the so-called INF projects.
This report tries to catalogue the INF projects in Germany's research landscape.
Building on that and on the workshop held at the de-RSE26 conference we collect lessons learned across INF projects and distill the discussion into guides that other INF projects (or similar structures) can follow.

\section{History of the INF projects}
The beginnings of specialised, shared positions within CRCs can be traced back to the foundation of the CRC 299 which was established in 1997 to survey the use of land under various characteristics. Being interdisciplinary and already aiming for the maximum runtime of 12 years it was apparent from the beginning \cite{mueckschel2006} that a careful
management of data has to happen in order to be useful for all associated disciplines.
Around the same time, various incidents of scientific misconduct prompted the German Research Foundation (GRF, DFG) to highlight in their
memorandum on ``Safeguarding Good Scientific Practice'' \cite{DFG1998}
the importance of research data as ``primary data''  and has put out the recommendation to take special care
of those data at the host institution.
Combining both requirements more CRCs started to include personnel that handle data management.
A first workshop \cite{SFB299-2004, SFB299-2005} on comparing lessons learned was organised by CRC 299 and CRC 552 in 2004 in Gießen.
The respective follow-up workshop was organised again by CRC 299 and CRC 552 in 2007 \cite{SFB299-2007}.
Around the same time the process of having dedicated projects for the information infrastructure was formalised(based on the precursor work of CRC 299 and CRC 552)
and since 2007 INF projects can be applied for in CRCs at the GRF with the first projects starting around 2009 \cite{Forschungsdateninfrastruktur2013}.
Moving on, the regional focus shifts to cologne.
In 2009 a workshop was organised by TRR32 and CRC806
bringing together almost 80 participants from a diverse set of disciplines which published their experiences in a
Post-Proceedings volume \cite{curdt787}.
The next meet-up was organised by the project Radieschen
``Rahmenbedingungen einer disziplinübergreifenden Forschungsdateninfrastruktur''
in Göttingen \cite{Forschungsdateninfrastruktur2013}. Among the participants they observed,
that most INF projects were with central institutions like libraries and compute centers and maintained central
research data management infrastructures for research collaboration. And, again, acceptance of the INF projects was brought up.\\
Between 2013 and 2015 the number of funded INF projects increased and hence the funding line became more established \cite{Engelhardt2013Forschungsdatenmanagement, bastian_2025}

Another follow-up workshop was held in Göttingen in 2018 \cite{Roertgen2019}.
Around 2022 the DFG started to hands out guidelines for the INF projects \cite{bastian_2025}.
Another meet-up that should belong in this list was organised by the DINI/nestor AG Forschungsdaten in 2020\cite{Rizzolli_2021},
that discussed how much the tasks of an INF project align with the tasks of data stewards.
Again in Cologne, another workshop was hosted \cite{bastian_2025}. A notable insight is, that no particpant has been to any
of the previous workshops, and it was noted, that following the rules is still a better motivator for following data
management guidelines, than any intrinsic motivation from the scientists.
In March 2026 we held a small workshop at de-RSE26. Notable is here, that as many Research Software Engineers have been present
as Research Data Managers.

Already early on \cite{mueckschel2006} it was noted that central challenges of a data management solution will be
\begin{itemize}
 \item personnel fluctuations over the runtime of a CRC,
 \item safe storage for the lifetime of a CRC,
 \item adaptation to the scientists questions over the evolution of their research
 \item extensibility and openness of any system,
 \item acceptance by the users
\end{itemize}
of any measure of the information system administered by the data stewards.
These are topics that persist until today,
and we will try to find out which answers the current generation of INF projects finds.


\section{INF Projects in Germany}
In this section we describe the various INF projects in Germany, how they see themselves and which duties they perform.
In this paper we exclusively refer to these projects as INF projects, but within the resepective CRCs, also other acronyms are used.

\subsection{Project Z03 --- Sustainable code development and optimization on modern supercomputing architectures, \href{https://gepris.dfg.de/gepris/projekt/258499086?language=en}{SFB1170 - Topological and Correlated Electronics at Surfaces and Interfaces, ``ToCoTronics''}, Universität Würzburg}
Established in 2015 as a support unit for the theoretically working projects of the CRC the focus was on numerical methods for Quantum Monte Carlo methods, functional renormalization group techniques and non-equilibrium Green's function techniques. These methods are usually used on
Germany's HPC systems, therefore the tight integration between software development, deep knowledge about the methods used in the CRCs science and a knowledge about the targeted platforms became paramount to this INF project staffed by one person.
To facilitate the spread of knowledge regular teaching opportunities were organised based on the materials of \href{https://carpentries.org/}{the Carpentries} in collaboration with the local biology department.
From the very beginning also the maintenance of the local IT infrastructure of the physics department was part of the project which included
storage systems, desktops, as well as an HPC system.
This integration of maintaining the IT infrastructure while at the same time focussing on sustainable development
facilitated the roll-out of a local gitlab instance for the entire department which has grown by 2026 to $\approx 1000$ users from all
over the world. Networking opportunities were used by participating in the RSE movement, which included the formation of the de-RSE society in 2018. This broad network enabled the project to write the data management strategy of the \href{https://www.ctdqmat.de/}{EXC2147} which was successfully reviewed in 2019.
Now also data management questions were part of this project.
The continued usefulness of this project was acknowledged by the DFG since a second funding round of the SFB1170 was granted also in 2019.
2020 was used for the formation of a local chapter ``Wü-RSE'' to provide a place
locally for like-minded people. The project was extended in 2023 for the last funding period. A highlight for this project's visibility included the hosting of de-RSE24, bringing the community of RSEs to Würzburg. Embedded into the community of RSEs, this project contributed
to numerous position papers over the years. FIXME:cite them\\
This project spans the breadth of digital tools in science from a single machine up to HPC clusters, from data generation, over data analysis upto data storage, collaborating with single scientists as well as larger research groups. This required a flexibility in the tooling, which obviously used git, but led to the use of C++, Python, Mathematica, Fortran, and Julia for programming. The CRC1170 and EXC2147 are supporters of the FairMat NFDI consortium, so this data platform is used in the department.
FIXME: SHORTEN....

\subsection{Project INF --- Scientific services and data management - \href{https://www.neglab.de/}{CRC 1629 Negation in Language and Beyond} - Frankfurt am Main, Tübingen, Göttingen}
CRC~1629 \emph{Negation in Language and Beyond} (NegLaB) is a Collaborative Research Center that investigates negation as a linguistic and cognitive phenomenon from different empirical and theoretical perspectives. Its projects span a wide range of disciplines, including theoretical linguistics, corpus-based research, and experimental work. As a result, CRC~1629 generates very different kinds of data through diverse methodologies and workflows. Within this setting, our project plays a central role by providing research data management, consulting, and teaching. This is particularly important in a humanities CRC, where researchers often work with sensitive data and are usually not formally trained in areas such as data management, data analysis, statistics, or experimental design.

Our project is interdisciplinary, bringing together computer scientists, theoretical linguists, and cognitive scientists. This is both its major strength and one of its main challenges. The challenge is further amplified by the fact that NegLaB’s scientific questions explicitly depend on cross-linguistic and cross-method integration. As a result, infrastructure decisions affect not only individual projects, but also how easily results can be compared and integrated across the CRC. In particular, it has been challenging to establish data-management practices such as naming conventions, data organization, and long-term data stewardship. To address these challenges, the project develops training workshops that are useful for researchers with very different backgrounds and needs.

An important part of this work has also been the collaboration with the library services of Goethe University Frankfurt. Together, we have worked on developing a CRC-wide data management protocol, as well as appropriate data management plans for each project using the Research data management organiser (RDMO). This collaboration has also helped increase the visibility of INF's research data management initiatives, as these have been presented and discussed in events organized by the library services.
Beyond our host institution we are in contact whith \href{https://www.uni-marburg.de/en/hefdi/about-hefdi}{HefDi}
which is a Hessen-wide network that offers courses and resources regarding data management, as well as with other Linguistics SFBs.

Our biggest takeaway is how important an INF project is in a humanities CRC: it helps ensure that sensitive and valuable data are handled appropriately, while at the same time giving researchers access to forms of training they would otherwise rarely receive. We have also learned how productive collaboration between computer scientists and linguists can be, as linguists provide fine-grained and theoretically grounded insights into language, while computer scientists contribute more efficient implementations and novel methodological approaches for research.


\subsection{Project INF - Sustainable Research Data Management, \href{https://gepris.dfg.de/gepris/projekt/417590517?language=en}{SFB1415 -  Chemistry of Synthetic Two-Dimensional Materials}, Technische Universität Dresden}
The CRC\,1415 "Chemistry of Synthetic Two-Dimensional Materials" established 2020 at TU Dresden operates at the intersection of chemistry, materials science, and physics
 aiming for controlled "bottom-up" synthesis and development of novel synthetic two-dimensional materials with high structural definition at the atomic or molecular level.
Since 2024 in the start of the 2nd funding period, the INF project accompanied the central project to develop and implement sustainable research data management (RDM)
infrastructure to support reproducibility and collaboration across 25 scientific projects and multiple partner institutions (Leibniz/Helmholtz) fostering harmonized data management strategies.
The primary focus has been on RDM—installing and maintaining scalable infrastructure
(e.g., local \href{https://nomad-lab.eu/nomad-lab/nomad-oasis.html}{NOMAD Oasis} of FAIRmats \href{https://nomad-lab.eu/nomad-lab/nomad.html}{NOMAD} platform),
integrating electronic lab notebooks (e.g., \href{https://www.elabftw.net/}{eLabFTW}),
and providing regular training workshops for doctoral researchers and project leaders on FAIR data principles and infrastructure use.
Research Software Engineering (RSE) tasks, such as developing data pipelines, writing file parsers for proprietary formats as well as automatic data processing, complement these efforts.
The CRC\,1415 and the INF project is actively embedded in national and international networks,
including the National Research Data Infrastructure (NFDI, particularly FAIRmat and NFDI4Chem), DINI, nestor, and the RDA Germany,
which collectively facilitate knowledge exchange and the adoption of RDM best practices.
Locally, the INF project has strong institutional connections with TU Dresden’s CIDS ZIH (Center for Information Services and High Performance Computing)
and Service Center Research Data, and the Saxon State and University Library Dresden (SLUB)
combining technical, infrastructural, and domain-specific knowledge.
A significant challenge has been the consistent implementation of FAIR data principles and harmonized metadata standards across the CRC’s heterogeneous experimental and theoretical projects,
as well as the unified rollout of the electronic lab notebooks eLabFTW and local data sharing platform NOMAD Oasis, due to diverse workflows and technical requirements. 
In contrast, open-access paper publishing have been widely implemented and open-access data publishing via \href{https://zenodo.org/communities/crc1415/about}{Zenodo}, \href{https://opara.zih.tu-dresden.de/home}{OPARA}, and \href{https://rodare.hzdr.de/communities/crc1415}{RODARE} grows in adoption,
directly supporting the scientific outreach.
The experience in CRC\,1415 underscores that technical infrastructure alone is insufficient; sustained training, community engagement, and the early embedding of RDM and RSE expertise within the CRC’s data lifecycle are essential for the successful and sustainable adoption of open science practices and FAIR research activities.

\subsection{Project INF --- Collaborative and Sustainable Research Data Management, Analysis and Visualisation, SFB1340 -- Matrix in Vision, Berlin} \label{sec:sfb1340}
The INF project of CRC 1340 “Matrix in Vision” originated in medical informatics and radiology at Charité – Universitätsmedizin Berlin and the Berlin Institute of Health (BIH). It was conceived as a central service project for a data-intensive CRC aiming to understand how pathological changes in the extracellular matrix can be visualised in vivo, particularly in inflammatory disease settings, by combining biophysical and molecular imaging approaches. Within this context, the INF project provided a shared digital backbone for collaboration, data integration, research data management (RDM), secure data exchange, and reusable analysis workflows, especially for multimodal imaging, experimental data, and clinical data. In practice, the project focused mainly on RDM, complemented by infrastructure and support tasks. A central outcome was the self-hosted Matrix-COOP platform at Charité, based on Nextcloud, with about 100 TB managed storage and more than 100 active users. It served as a shared space for secure data exchange, templates, reporting, onboarding information, and a CRC-wide knowledge base. The project also supported data management plans, data protection and governance coordination, and access to centralized data storage, archiving and HPC resources. A particularly visible example was the integration of about 2,600 electronic health records with myocardial MRI data into a dedicated data warehouse with self-service query functionality and integrated data provenance tracking \cite{johns_prov_tracking_2025}. Another important line of work concerned the secure exchange of imaging data and the setup of interfaces and transfer workflows for large multimodal datasets across the CRC. The project was closely connected to local and national research data infrastructures, especially NFDI4Health \cite{fluck_bausteine_2021} and the Medical Informatics Initiative \cite{prasser_difuture_2018}. Locally, it strengthened connections between medical informatics, radiology, the IT department, and CRC subprojects by embedding service development directly into day-to-day research workflows. A key lesson was that a single unified RDM platform for all groups was not realistic because requirements, data types, and established practices differed too much between projects and according processes were already established in the first funding phase, in which no INF project was included. The project therefore shifted from a one-size-fits-all model to a hybrid, federated approach combining shared infrastructure with project-specific support. Easier to adopt were pragmatic services with immediate benefit, especially the cooperation platform, shared documentation, and targeted support for concrete integration needs. Overall, the project showed that sustainable INF work in a CRC requires both technical infrastructure and continuous hands-on support close to the researchers. Its usefulness was reflected in the successful evaluation of the CRC for a third funding period, with relevant parts of INF being carried forward as part of a new subproject.

\subsection{Project INF --- Research data management and collaboration platform, TRR418 -- Foundations of Circadian Medicine, Berlin/Lübeck/Munich}
The INF project of TRR 418 “Foundations of Circadian Medicine” combines expertise in medical informatics, health data privacy, and bioinformatics at Charité – Universitätsmedizin Berlin and the Berlin Institute of Health (BIH). It supports a transregional consortium investigating how circadian clocks shape physiology and disease and how this knowledge can be translated into medical research and practice. Because the TRR combines preclinical and clinical work across several sites, the INF project was conceived as an enabling structure for coordination, data stewardship, and cross-project reuse. Its main focus lies on a platform for broadly sharing circadian profiles and reproducible computational workflows, complemented by privacy-related guidance and training \cite{baum_data_circadian_2023}. A defining feature of the TRR is its reliance on longitudinal and rhythmic data, ranging from high-resolution sensor and wearable data to clinical, phenotypic, and omics datasets. The INF project therefore aims to create an environment in which such heterogeneous data can be documented, exchanged, processed, and explored in a consistent way. Planned components include a collaboration platform and services for FAIR data management with a particular emphasis on temporal data and tools for exploring circadian and omics data across projects. Analogously to the INF project of \hyperref[sec:sfb1340]{CRC 1340}, this project is embedded in institutional and national infrastructures. An expected challenge is the adoption of common data representations, ontologies and workflows across projects that differ widely in modality, temporal resolution, and regulatory context. Easier to adopt will likely be services with immediate practical benefit, such as shared collaboration spaces, and support for reproducible data and code handling. Since the project started only in October 2025, its contribution so far is mainly the coordinated establishment of shared structures; the central lesson at this stage is that, in a consortium centered on circadian and longitudinal data, INF work must integrate RDM, RSE, privacy, and domain-specific analytics from the outset rather than treating them as separate support functions.

\subsection{Project INF --- Research Data Management and Research Software Engineering, \href{https://gepris.dfg.de/gepris/projekt/327154368}{SFB1313 - Interface-Driven Multi-Field Processes in Porous Media – Flow, Transport and Deformation}, University of Stuttgart}
CRC 1313 was established in 2018 to study interface-driven processes in porous media.
To manage the data and code generated by these studies, a task force named "Software and Data" was formed to develop a long-term strategy. This work led to the creation of the INF project, which began in 2022 during the second funding phase. The project focuses on RDM and RSE to ensure that data and software outputs follow the FAIR principles. By reviewing datasets and providing support for software development, the project helps researchers produce high-quality, reusable, and well-tested code. These efforts are supported by training programs and outreach activities.

The project is well-integrated into both national and local networks, including NFDI4Ing, de-RSE, and Data Stewardship goes Germany.
At the local level, INF has built strong connections with the \href{https://www.simtech.uni-stuttgart.de/communication/events/workshops/sigdius-seminars/}{SIGDIUS} network and the central IT service facility of the University of Stuttgart. While some practices were adopted quickly, such as using quality badges and publishing datasets in the DaRUS repository, others proved more difficult. For example, implementing ELNs (e.g. \href{https://www.elabftw.net/}{eLabFTW}) has been challenging because it requires a significant change in daily research habits. A major lesson learned is that technical infrastructure is not just about tools, but also about people. To convince researchers to adopt new practices while facing tight deadlines, the project must demonstrate clear, short-term benefits. Ultimately, success depends on effective communication and building collaborations that show how better data and software management improves research efficiency.

\subsection{Project INF - User-oriented research infrastructure assisting linguistic data collection and (re-)use, \href{https://www.uni-bielefeld.de/sfb/sfb1646/}{SFB1646 -  Linguistic Creativity in Communication}, Universität Bielefeld}
CRC 1646 was established in 2024 at Bielefeld University. The research is an interdisciplinary collaboration involving multiple fields, such as theoretical and experimental linguistics, computational linguistics, psycholinguistics, and others. It investigates the core properties of linguistic creativity, develops empirical measures, and analyzes underlying linguistic, cognitive, and social mechanisms. The research is structured in three phases: the first formulates initial hypotheses, the second examines principles and their interactions in creative language use, and the third integrates the findings into a comprehensive theory of linguistic creativity. The projects conducted within this CRC not only require handling large (e.g., video data) and diverse (e.g., video, audio, picture, text, fMRI, etc.) datasets but also involves the collection of personal data, introducing an additional layer of sensitivity to data management. Some of the projects include experimental studies that gather demographic information, video recordings, or medical history data of their participants, all of which necessitate carefully designed consent procedures. Researchers must ensure both the legal \cite{GDPR2016} validity of consent forms and the secure management of sensitive data, while still enabling data sharing in accordance with FAIR principles \cite{WilkinsonDumontier2016}.

To address these challenges, the INF project aims to establish a dedicated infrastructure comprising two main components.
First, a \textit{wizard} platform that supports researchers during the data collection phase by automatically generating customized, legally compliant consent forms tailored to specific project requirements.
Second, a data management platform that enables the integration and handling of collected (meta)data. Within this platform (in the second phase), data will be systematically linked to each participant's consent, allowing researchers to clearly determine the scope and limitations of data usage. In this phase, INF focuses on three main tasks: RSE for the development of the wizard, providing RDM support, and offering programming, methodological, and statistical support.

\textbf{I.} Our wizard \textit{RUDI} (Research-centred User-oriented Data Infrastructure) is being developed as a dynamic, modular system, grounded in a careful analysis of domain-specific requirements and designed for future adoption by other disciplines.
The infrastructure is based on ontological data models, enabling different layers and components of the platform to automatically adapt to new values while maintaining consistency and integration.
The considered model enhances the next phase by improving project collaboration efficiency, and aligning with strategies that tailor data management to the specific needs of consented data \cite{Mohammadietal2026}. %, thereby facilitating effective collaboration and reducing the time required for data management tasks.
\textbf{II.} Given the large volume and heterogeneity of data within the CRC, we provide a range of RDM-related workshops, including training on GitLab (the recommended repository at Bielefeld University), tidy data principles, and FAIR data practices. These activities are particularly aimed at preparing early-career researchers, such as PhD students, for data management responsibilities. In addition, we offer comprehensive \href{https://crc1646.pages.ub.uni-bielefeld.de/inf/website/wiki/site/index.html}{guidelines} and a \href{https://crc1646.pages.ub.uni-bielefeld.de/inf/website/faq.html}{FAQ section} on our website, enabling members to access relevant information as needed. In accordance with university policies, INF maintains close collaboration with \href{https://www.uni-bielefeld.de/ub/digital/forschungsdaten/kompetenzzentrum/}{Competence Center of Research Data} and \href{https://www.uni-bielefeld.de/einrichtungen/bits/ueber-uns/index.xml}{Bielefeld IT service (BITS)}. 
\textbf{III.} Finally, INF provides training in statistical methodologies through workshops and hands-on sessions. We also offer workshops and consultation on programming and annotation tools to support researchers in their technical work. In addition, our legal colleagues in INF offer guidance on the legal aspects of consent forms, highlighting key requirements and best practices through dedicated workshops and consultation sessions.

\subsection{Project INF - Infrastructure, TRR265 - Loss and regaining control in addictions, TUD Dresden University of Technology}

The INF project provides a comprehensive IT infrastructure tailored to the needs of the researchers working in the domain of addiction research (primarily alcohol and tobacco) within the clinical psychology. It builds on and extends the infrastructure created for the prior SFB940 ``Volition and Cognitive Control'' as well as for the CRC205 ``The Adrenal: Central Relay in Health And Disease''. The goal of INF is to establish a unified system supporting every step of the clinical research. This begins with the recruitment of participants, their screening and pseudonymization, the management of the studies, the conduction of digital behavioral tests, as well as the processing and long-term archival of fMRI data.

The project fostered the synergy between the Center for Information Services \& High Performance Computing (ZIH) of the TU Dresden, Dresden's University Clinic and the Faculty of Psychology, especially its Neuroimaging Center (NIC). It also demonstrated the need for a clear documentation and communication of the agreed procedures and the importance of strong quality assurance procedures at every step of the workflows involving human interaction.


\subsection{Project INF - Infrastructure, TRR393 - Trajectories of Affective Disorders: Cognitive-Emotional Mechanisms of Symptom Change, TUD Dresden University of Technology}

TRR393 is a Collaborative Research Centre (SFB/TRR 393) focused on Affective Disorders (AD) — comprising major depressive disorder and bipolar disorder. Its goal is to identify trajectories of recurrences and remissions in AD, determine cognitive-emotional mechanisms and neurobiological correlates of acute symptom changes, and probe mechanism-based interventions.

The INF project establishes a unified information infrastructure for the CRC/TRR 393, designed to support the research of individual projects, foster cross-project collaboration, and ensure adherence to FAIR data management principles. At its core, the project pursues three key objectives: (i) providing a professional system for research data capture, storage, access, and information management; (ii) implementing comprehensive data quality control procedures; and (iii) delivering training for standardized clinical assessments across all participating sites.
The infrastructure encompasses a broad range of components that together form a cohesive research ecosystem. These include centralized study management in REDCap, GDPR-compliant handling and pseudonymization of personal data in Participant Management System (PAMS), centralized quality control and archiving of MRI data, seamless integration of the MRI devices at UMR and UMS into the archiving system, and the import of existing datasets from BipoLife and FOR 2107. Developed and operated in close cooperation with all CRC/TRR 393 researchers, the system evolves continuously to meet emerging project requirements — accommodating new data types along with their metadata — and is complemented by ongoing user support and training in the IT environment.
Building on methods and tools successfully established in CRC 940 and CRC/TRR 265, the INF project further benefits from ZIH's commitment to sustainable long-term data archiving, ensuring that the scientific output of CRC/TRR 393 remains accessible, secure, and reusable for years to come.



\subsection{Project Z0X - project title, SFBZZZZ - SFB/TRR title, cities}
- what domain are you from, what was the project idea you were based on?\\
- What are the tasks you did focus on, RDM vs RSE vs Teaching vs. sth else?\\
- (For the Definitions: RSE, see here: \url{https://de-rse.org/learn-and-teach/}, \url{https://de-rse.org/de/} , RDM, see here: \url{https://forschungsdaten.info/} )\\
- Which networks are you part of?\\
- Which connections did you locally form at your site in your institution/university?\\
- Give an example of a thing which you consider useful but has been very hard to adopt.\\
- Give examples of things that were easily adopted in your CRC.\\
- What have you learned in this project?\\
Hint: Limit this section to around a quarter of a page


\subsection{Project Z0X - project title, SFBZZZZ - SFB/TRR title, cities}
- what domain are you from, what was the project idea you were based on?\\
- What are the tasks you did focus on, RDM vs RSE vs Teaching vs. sth else?\\
- (For the Definitions: RSE, see here: \url{https://de-rse.org/learn-and-teach/}, \url{https://de-rse.org/de/} , RDM, see here: \url{https://forschungsdaten.info/} )\\
- Which networks are you part of?\\
- Which connections did you locally form at your site in your institution/university?\\
- Give an example of a thing which you consider useful but has been very hard to adopt.\\
- Give examples of things that were easily adopted in your CRC.\\
- What have you learned in this project?\\
Hint: Limit this section to around a quarter of a page

\section{Lessons learned - Guides}
In this section we summarise the lessons we have learned together in this workshop.

\subsection{Outreach, Collaboration and Networking}
INF projects are often one of the cross-cutting components within a CRC, and such they play an important part in the collaboration and networking within the project.
Also, because partners working in INF projects necessarily have to have a higher-level overview of work within the entire CRC, they can be valuable partners or even facilitators of outreach of the project.
This section discusses experiences in this important social and cross-disciplinary component of an INF project.

\todo{maybe interesting: \cite{Hausen_Rosenberg_Trautwein-Bruns_Schwarz_2020}}
We acknowledge that some of the experiences likly only apply to certain categories of domains.
However, we also expect that often these are wide enough, e. g. “the natural sciences”, to be interested to the reader.
It would be interesting to look at the breath of applicability of experiences, but we do not have enough data to make any such claim, so most of the following has more anectodal character.
Nevertheless, participants expressed that even sharing those was valuable for them and likely their projects.

We start with the usual suspects and their perceived effectiveness, rooted in project management.
Most projects hold annual status and PhD candidate meetings.
Both are seen as somewhat usefull, but in the end that frequency is seen as too low.
In contrast, a quarterly frequency is experienced as much more effective.
Three months is seen as not too often to get in the way of “regular work”, but frequently enough for most projects to provide early feedback and incentive for a timely implementation of agreed-upon targets.
Most importantly, this frequency is not only seen as more effective than yearly by the INF teams, but also by the researchers, increasing their engagement with the INF projects.
Also connected to update intervals, an open development (as opposed to a closed one) eased outreach for multiple INF projects, as researchers felt more connected to the work within the INF project.

Another topic besides meeting frequency is the choice of contact channels.
The general agreement within the workshop participants was to disagree on the specific prefered channels, but everyone agreed that the preferred choice is effectively dictated by the channels the researchers within the CRC prefer.
Especially, social media plattforms seems to work well if the targeted research commmunities already use them.
Thus, one advise from the workshop is: Investigate the prefered communication and outreach channels of your community and use those.
While creating networks between non-INF projects was often experienced as easier than between INF and non-INF projects, serving and then re-using their communication channels is often very helpful.

Generally, meeting formats with a lower (near-work) level seem to work far better than higher-level formats.
This includes offerings as part of PhD-schools, courses, as well as internal project pairings.
The closer those are to what researchers perceive as “the actual work”, the more successful they are.
One of the prerequisites for success, also more generally, of these efforts seems to be an explicit backing of such low-level offerings and pairings by the PIs of the project.
A challenge that stems from this is that new topics within digital research, e. g., tools or methods, can sometimes be relatively easy to introduce on a low-level, but may require substancial persuasion on the level of research-PIs.
One take-away point here is that this PI-buy-in can not only be helpful but often is inevitable.

Not surprisingly, general ideas that are close to the Open Science community, and thus to the RDM and RSE communities as well, like FAIR\todo{cite here if not already used before}, can be hard to convey within some CRCs, especially when the served communities do not put a high value on those themselves.
There, an INF project often not only has the role as RDM/RSE-enabling unit, but also as Open Science evangelist.

\subsection{How do we educate scientists and ourselves}
Education and training for scientists and research data managers is a multifaceted challenge that goes beyond simply providing access to courses.
We acknowledge that an important task of INF projects is the provision of teaching materials, but the training of the INF project itself is equally important.
While many researchers are capable of identifying and enrolling in online training opportunities independently, a recurring difficulty lies in finding relevant and high-quality training that aligns with their specific needs and contexts.
This highlights the importance of not only curating suitable offerings but also guiding researchers toward them.
A key factor in effective training is the careful tailoring of content to the target audience.
Generic training formats often fall short, particularly when domain-specific tools and practices are involved.
Integrating such specialized tools into training programs remains a significant challenge, yet it is essential for ensuring applicability in real-world research settings.
Without this alignment, there is a risk of participants loosing interest and adopting suboptimal or “wrong” solutions early on —-- leading to a kind of ``Good-Enough'' problem, where initial partial success obscures underlying issues and may even have long-term negative consequences.

To address these challenges, training efforts should be complemented by consulting services that engage directly with researchers’ specific use cases.
This allows for more targeted support and helps bridge the gap between theoretical knowledge and practical implementation.
In many cases, basic training is already provided by infrastructure (INF) teams themselves, but there is clear value in expanding beyond these internal efforts.
Leveraging external training opportunities -- such as those offered externally to the CRC or broader initiatives -- can help diversify and strengthen the available training landscape.

Another important dimension is the relationship between different communities, particularly between INF projects and domain researchers.
In some fields, formal certification processes are difficult to establish, partly due to mutual uncertainties and a lack of shared understanding.
This underscores the need for improved communication and information exchange between these communities.

Training is not only relevant for researchers but also for RSEs, Data Stewarts and other support roles, whose professional development is equally important for sustaining high-quality research environments.
In some cases, dedicated positions exist with explicit responsibility for training, extending beyond purely infrastructure-related topics.
However, these roles are often developed “on the job,” and a recurring issue is the lack of recognition through traditional academic outputs such as publications.

Overall, strengthening education and training in this space requires a combination of better curation, audience-specific design, integration of domain tools, and closer alignment with real research practices. Equally important are structural improvements, such as fostering cross-community exchange and recognizing training-related work as a valuable professional contribution.


\subsection{Reproducibility, Quality and Validation of the artifacts we are creating}
Having received proper training, how do we validate the artifacts that are created during our work?\\
Possible keywords: long term archiving measures, containerization, research and software reproducibility, Verification and Validation (e.g. Performance), Software quality, badges,\\

The workshop discussions established that while automated tools are rapidly advancing, human expertise remains central to artifact validation. In fields dealing with sensitive clinical or animal data, researchers are increasingly adopting a "human-in-the-loop" approach. Within this framework, machine learning models and locally-hosted or privacy preserving Large Language Models (LLMs) are utilized to suggest metadata or classify complex datasets, which are then verified through expert visualization and manual checks. This method ensures that data privacy is maintained while leveraging AI to handle the increasing volume of research outputs.

Technical validation becomes more complex as software grows in scale, requiring sophisticated containerization and environment management. Participants identified tools such as Pixi and the MaRDI packaging system as essential for managing dependencies and namespaces in large-scale projects. For dynamic datasets that change over time, static identifiers alone are often insufficient to guarantee reproducibility. Instead, modular Jupyter notebooks and data chunking strategies are used to document specific states and edge cases, ensuring results can be reconstructed even when the underlying data is not static. Furthermore, implementing CI/CD pipelines and automated testing allows for the continuous monitoring of software quality, specifically checking for breaking changes or inconsistencies caused by random seeds and stochastic processes.

Finally, the long-term quality of artifacts depends on the rigorous use of Electronic Lab Notebooks (ELNs) like Chemotion, eLabFTW, and Sciformation to capture  experimental procedures and data provenance at the point of origin. Integrating ELNs with established repositories like Dataverse/DaRUS enables structured, metadata-rich publication of experimental data, while version control platforms like GitHub support the parallel management of research software. Despite these tools, the workshop highlighted a persistent gap in the quality of researcher-provided metadata, pointing to the need for better training and workflow integration to ensure artifacts remain verifiable and reusable over the long term.


\subsection{Workflows and pipelines}
Having created pieces of software, we need to integrate them into longer stages of pipelines

Possible keywords: RDM, workflow(e.g. from NFDI), best practices, AI practices and research workflows\\

A big benefit are workflow platforms. Examples of this are Galaxy, Taverna, and workflowhub. The use of standardised languages like CWL would be beneficial.
Examples of these languages are Mardiflow, snakemake and nextflow. Since github and gitlab pipelines offer similar features, they are also useful.
But with this there come downsides:
So this requires training. Since you have to detail your steps in a standardized language you also have to learn it. So this requires Coding practise.
THe adaptation of pipelines and version changes can bring the well-known problem of Maintenance.
Workflows have to be permitted in the respective domain, So it can be, that the use of a workflow requires approvals, e.g. by an ethics board or similar.
Workflow languages currently definitively have a low reach into the user base
Findability of hte respective resposible users is a problem.
It can be difficult to get the permission to use input data
Re-Use and Maintenence are an issue.
There is a lack of awareness for the benefits of automation, that often leads to people not doing the first steps.
Possible solutions are the use of an RAG supported Chat-Agent or scientists, or an easy/good Documentation.

\section{Summary}
I would like to end this on something that binds the INF projects together irrespective of the domain they are in.
Any other ideas for a summary?

\begin{acknowledge}
We acknowledge the hospitality of the university of Stuttgart for hosting this workshop.
FG acknowledges funding from the Deutsche Forschungsgemeinschaft (DFG, German Research Foundation) through the SFB 1170 “Tocotronics”, project Z03 - project number \geprislink{258499086}.
RD thanks the DFG for its support through the SFB 1415 "Chemistry of Synthetic Two-Dimensional Materials", project INF - project number \geprislink{417590517}.
AA acknowledges funding from the German Research Foundation (DFG) through the TRR 265 "Losing and Regaining Control over Drug Intake: From Trajectories to Mechanisms to Interventions", project  Z/INF - project number \geprislink{402170461}.
LS acknowledges funding from the German Research Foundation (DFG) through the TRR 393 "Trajectories of Affectie Disorders: Cognitive-Emotional Mechanisms of Symptom Change", project  Z/INF - project number \geprislink{521379614}.
HO acknowledges funding from the Deutsche Forschungsgemeinschaft (DFG, German Research Foundation) through the SFB 1313 “Interface-Driven Multi-Field Processes in Porous Media – Flow, Transport and Deformation”, project INF - project number \geprislink{327154368}.
MJ, TP, FP acknowledge funding from the DFG through the SFB 1340 "Matrix in Vision", project INF - project number \geprislink{372486779}.
MG, MP, LFB, AM, FP acknowledge funding from the DFG through the TRR 418 " Foundations of Circadian Medicine", project INF - project number \geprislink{541063275}

MM acknowledge the funding from DFG, SFB 1646/1 2024 -- \geprislink{512393437}, project INF. 

\end{acknowledge}

\begin{aiuse}
AI has been used for copyediting purposes.
\end{aiuse}

\begin{data}
\end{data}

\bibliographystyle{eceasst}
\bibliography{report}

\end{document}
